\section{Predictive Value of Sentiment}
\label{sec:predictive}

The preceding sections have demonstrated that the sentiment series is coherent, that it reflects historical episodes, that it differs systematically between the prepared and spontaneous sections of the press conference, and that it leads the policy rate around decision events. We now ask the operational question that follows from these descriptive results: \emph{does the textual signal predict the direction of the next ECB rate decision once we control for the rate history itself?} We address it in three steps. Section~\ref{sec:results_correlations} characterises the predictive correlation between sentiment features and rate changes across monetary regimes and lag horizons. Section~\ref{sec:results_prediction} fits four classifiers of increasing complexity---a logit baseline, a sentiment-only logit, an augmented logit, and a random-forest extension---to the meeting-level rate-decision direction. Section~\ref{sec:results_surprise} closes the chapter with a brief robustness check on communication surprise and market volatility.

\subsection{Cross-Era Predictive Correlations}
\label{sec:results_correlations}

Figure~\ref{fig:lag_profile} traces the Pearson correlation between three selected sentiment features and the contemporaneous \gls{mro} change across horizons of zero to fifteen meetings (changed meetings only, $n \approx 54$). Two patterns emerge. First, the correlation is already small but significant at lag~$0$ ($r = 0.31$, $p < 0.05$ for the FinBERT overall mean); it then rises to a pronounced peak at lag~$4$ ($r = 0.551$). Second, the profile is non-monotonic, with a secondary plateau through lags~$6$--$10$, indicating that sentiment carries predictive content at multiple horizons. The three features peak at different lags---FinBERT IS at lag~$2$ ($r = 0.352$), FinBERT overall at lag~$4$ ($r = 0.551$), and CentralBankRoBERTa IS at lag~$9$ ($r = 0.409$)---so the textual signal is not a single-horizon lead.

\begin{figure}[htbp]
    \centering
    \includegraphics[width=\linewidth]{%
        images/results/correlations/fig_lag_profile.pdf}
    \caption[Lag profile of sentiment features versus MRO changes]{%
        Pearson correlation between three lagged sentiment features and
        contemporaneous \gls{mro} rate changes, changed meetings only
        ($n \approx 54$). Stars mark the feature-specific peak.}
    \label{fig:lag_profile}
\end{figure}

Table~\ref{tab:summary} disaggregates the signal across monetary regimes. Three regularities emerge. First, correlations are positive in both the Pre-\gls{zlb} and Hiking regimes for almost every feature. Second, the
\gls{zlb} column is sign-heterogeneous: extremum features are negative or near zero while level-mean features at lag~$4$ are weakly positive---neither class tracks the shadow-rate accommodation reliably, for the reasons discussed in Section~\ref{sec:results_comovement}. Third, the four highest full-sample correlations belong to \emph{within-meeting maxima} at a two-meeting lag rather than means: the peak hawkish or dovish passage within an Introductory Statement carries more directional signal than the meeting average.

\begin{table}[htbp]
    \centering
    \caption[Top sentiment features by absolute predictive correlation]{%
        Top sentiment features ranked by absolute full-sample correlation
        (word limit $= 200$). \texttt{[t-k]}: lagged $k$ meetings;
        \texttt{FB} = FinBERT; \texttt{RB} = CentralBankRoBERTa;
        superscript $\max$ = within-meeting maximum score.
        Full: $\Delta\text{MRO}$, changed meetings, $n = 54$.
        \gls{zlb}: Wu--Xia shadow rate \cite{wu2016measuring}.
        Hiking $n = 13$; interpret with caution.}
    \label{tab:summary}
    \begin{tabular}{@{}lrrrr@{}}
        \toprule
        \textbf{Feature}              & \textbf{Pre-ZLB} & \textbf{ZLB}      & \textbf{Hiking} & \textbf{Full} \\ \midrule
        FB IS $\max$ [t$-$2]          & $+0.68$          & $-0.23$           & $+0.58$         & $+0.61$       \\
        RB IS$\cdot$EP $\max$ [t$-$2] & $+0.72$          & $\phantom{+}0.00$ & $-0.13$         & $+0.59$       \\
        FB $\max$ [t$-$2]             & $+0.64$          & $-0.21$           & $+0.58$         & $+0.58$       \\
        RB IS $\max$ [t$-$2]          & $+0.78$          & $-0.07$           & $-0.13$         & $+0.57$       \\
        FB QA$\cdot$MP $\mu$ [t$-$4]  & $+0.55$          & $+0.21$           & $+0.64$         & $+0.56$       \\
        FB MP $\mu$ [t$-$4]           & $+0.71$          & $+0.22$           & $+0.79$         & $+0.56$       \\
        FB $\mu$ [t$-$4]              & $+0.85$          & $+0.19$           & $+0.77$         & $+0.55$       \\
        FB QA $\mu$ [t$-$4]           & $+0.67$          & $+0.17$           & $+0.65$         & $+0.55$       \\
        RB $\max$ [t$-$5]             & $+0.69$          & $-0.06$           & $+0.38$         & $+0.54$       \\
        RB EP $\max$ [t$-$5]          & $+0.66$          & $-0.17$           & $+0.38$         & $+0.53$       \\
        \bottomrule
    \end{tabular}
\end{table}

The feature most consistent across all three eras is \textbf{FB $\mu$
    [t$-$4]} (Pre-\gls{zlb}: $+0.85$; \gls{zlb}: $+0.19$; Hiking: $+0.77$), with \textbf{FB MP $\mu$ [t$-$4]} following closely. Both are retained for the predictive model of Section~\ref{sec:results_prediction} on the basis of cross-era robustness. \textbf{FB IS $\max$ [t$-$2]} is retained as a short-horizon complement; its negative \gls{zlb} loading is handled by the regime-aware cross-validation in that section.

\subsection{Predicting the Direction of ECB Rate Decisions}
\label{sec:results_prediction}

The predictive task is operationalised as a three-class classification problem. Each press conference is assigned a label $y_t \in \{-1, 0, +1\}$ corresponding to the sign of the change in the \gls{mro} rate at that meeting: $-1$ if cut, $0$ if held, and $+1$ if hiked. The class distribution is heavily skewed towards \emph{hold} ($79.2\%$ of $280$ meetings), which sets a demanding naive baseline: a classifier that always predicts \emph{hold} achieves $79.2\%$ raw accuracy.

Four classifiers are compared. \textbf{M1 (Baseline)} is a logistic regression on the rate history only $\{\Delta\text{MRO}_{t-1},\, \text{MRO}_{t-1}\}$. \textbf{M2 (Sentiment)} is a logistic regression on four sentiment features---FinBERT \gls{is} mean at lag $2$, FinBERT overall mean at lags $2$ and $5$, and the FinBERT \gls{is} standard deviation---isolating the marginal content of the textual signal. \textbf{M3 (Augmented)} adds the M1 rate-history features to M2, testing whether sentiment adds information beyond the rate history. \textbf{M4 (Random Forest)} applies a 400-tree ensemble with bootstrapping and restricted depth ($\texttt{max\_depth}=4$) to the M3 feature set, relaxing the linearity assumption while guarding against over-fitting on the small sample of changed meetings ($n_{\text{changed}} \approx 58$). All models use class-balanced loss weights and are evaluated by $5$-fold time-series cross-validation, which respects the temporal ordering of meetings and never conditions on future data.

\begin{figure}[htbp]
    \centering
    \includegraphics[width=0.95\linewidth]{images/results/model/fig_performance.pdf}
    \caption[Cross-validated performance of the four predictive models]{Cross-validated performance of the four predictive models ($n = 275$--$278$). Solid bars: CV accuracy; outlined bars: AUC (one-vs-rest). Error bars show the standard deviation of accuracy across the five folds. The dashed line marks the naive baseline of always predicting \textit{hold}.}
    \label{fig:performance}
\end{figure}

Figure~\ref{fig:performance} reports accuracy and AUC for all four models. Three observations follow. First, M2's raw accuracy ($0.46 \pm 0.18$) is well below the naive baseline. This is the expected outcome of a class-balanced classifier on a heavily imbalanced sample: M2 actively predicts the minority classes (\textit{hike} and \textit{cut}) and is penalised on accuracy whenever those predictions are incorrect, even when the model successfully identifies the direction. Raw accuracy therefore understates M2's value, and the primary metric for comparison is the AUC. Second, the AUC ordering is strict: M1 ($0.680$) $<$ M2 ($0.706$) $<$ M3 ($0.719$) $<$ M4 ($0.741$). Each successive model adds a meaningful increment: the largest single step arises from adding sentiment to the rate history (M1 $\rightarrow$ M3, $+0.039$), while relaxing the linearity assumption adds a further $+0.022$ (M3 $\rightarrow$ M4). Third, the gap between M2 and M3 is small: the sentiment-only model captures most of the uplift over the rate-history baseline, confirming that the textual features rather than the rate-history features drive the incremental predictability.

To compare the two best models in detail, Figure~\ref{fig:comparison_m3_m4} places the confusion matrices of M3 and M4 side by side.

\begin{figure}[htbp]
    \centering
    \includegraphics[width=\linewidth]{images/results/model/fig_comparison_m3_m4.pdf}
    \caption[Confusion matrices of M3 and M4]{Confusion matrices of M3 (Augmented Logit, left) and M4 (Random Forest, right). Rows are true labels; columns are predicted labels. Per-class recall is reported in the subtitle of each panel.}
    \label{fig:comparison_m3_m4}
\end{figure}

The confusion matrices in Figure~\ref{fig:comparison_m3_m4} reveal a clear qualitative difference between the two models. M3 identifies $24$ of $29$ cuts (recall $83\%$) and $22$ of $28$ hikes ($79\%$), but only $127$ of $218$ holds ($58\%$): $48$ holds are labelled as cuts and $43$ as hikes. M4's full-sample confusion matrix is markedly better---Cut $97\%$, Hold $77\%$, Hike $89\%$---but these full-sample figures should be read with care. A bootstrap-aggregating ensemble fits its own training sample very closely, so the training-set confusion matrix overstates the out-of-sample precision; the honest comparison between M3 and M4 is the cross-validated AUC ($0.719$ versus $0.741$). What the side-by-side confusion matrices do illustrate reliably is the direction of improvement: M4 meaningfully reduces the false positives on the \textit{hold} class relative to M3, recovering meetings that the linear specification misclassifies as directional decisions. The substantive interpretation remains the same for both models: they are reliable for the directional question (\textit{is the next move up or down?}) but less precise about the timing within a sequence of hold meetings, exactly the pattern expected from a signal that leads the rate by several meetings.

The signed coefficients of M3 and the feature importances of M4 are reported separately in Figures~\ref{fig:coefficients} and~\ref{fig:rf_importances} to preserve the interpretive clarity of each representation.

\begin{figure}[htbp]
    \centering
    \includegraphics[width=\linewidth]{images/results/model/fig_coefficients.pdf}
    \caption[Standardised coefficients of M3 by class]{Standardised multinomial logistic coefficients of M3 (Augmented Logit) by target class (cut, hold, hike).}
    \label{fig:coefficients}
\end{figure}

\begin{figure}[htbp]
    \centering
    \includegraphics[width=0.80\linewidth]{images/results/model/fig_rf_importancies.pdf}
    \caption[M4 Random Forest feature importance]{Impurity-decrease feature importance of M4 (Random Forest). Features are labelled consistently with Figure~\ref{fig:coefficients}.}
    \label{fig:rf_importances}
\end{figure}

The two representations are mutually consistent and reinforce the same economic narrative. In M3, FinBERT overall mean at lag $2$ loads positively on \textit{hike} and negatively on \textit{cut}: hawkish language two meetings ago is the strongest short-horizon directional predictor. FinBERT overall mean at lag $5$ carries the same sign pattern at a longer horizon, consistent with the six-meeting telegraphing window identified in Section~\ref{sec:event_study}. The lagged \gls{mro} rate loads positively on \textit{hold}, capturing the low sequential reversal rate of the \gls{ecb}. The standard deviation of \gls{is} sentiment loads negatively on \textit{hold}: internally volatile communication is a marker of a meeting at which an actual decision is forthcoming. In M4, FinBERT overall mean at lag $2$ is the single highest-importance feature (impurity decrease $\approx 0.20$), with FinBERT \gls{is} mean at lag $2$ and $\Delta\text{MRO}_{t-1}$ close behind. The Random Forest thus assigns the greatest discriminating power to the same short-horizon textual channel that the logistic specification identifies, rather than uncovering an entirely different non-linear signal. The $+0.022$ AUC gain of M4 over M3 is therefore best interpreted as a quantitative tightening of a pattern already present in M3, not as evidence of a qualitatively different informational source.

\subsection{Communication Surprise and Market Volatility}
\label{sec:results_surprise}

The predictive analysis has so far focused on the level of sentiment and its lagged correlations with policy rates. A complementary perspective is to ask whether the \emph{change} in sentiment from one meeting to the next---the communication surprise---correlates with the contemporaneous market reaction during the press-conference window. Figure~\ref{fig:delta_sentiment} reports two such checks. The top row regresses the change in \gls{ois}-$1$Y on three lagged measures of $\Delta$sentiment; the bottom row regresses the within-meeting standard deviation of sentiment on the cross-maturity mean of $|\Delta\text{OIS}|$.

\begin{figure}[htbp]
    \centering
    \includegraphics[width=\linewidth]{images/results/model/fig_delta_sentiment.pdf}
    \caption[Communication surprise and market volatility]{Top row: change in sentiment versus the change in \gls{ois}-$1$Y during the press-conference window. Bottom row: within-meeting standard deviation of sentiment versus the cross-maturity mean of $|\Delta\text{OIS}|$. Points are colour-coded by era. Regression lines shown for visual reference.}
    \label{fig:delta_sentiment}
\end{figure}

The results echo Section~\ref{sec:divergence}: the directional correlations between communication surprise and the directional OIS move are statistically indistinguishable from zero (top row, all $|r| < 0.05$), but the standard deviation of sentiment is mildly negatively correlated with the absolute size of the OIS move ($r = -0.137$ for the CentralBankRoBERTa overall standard deviation, $p < 0.05$, bottom-right panel). The sign mirrors the communicative-cancellation finding: more internally variable communication is associated with smaller market reactions. This is consistent with a market that interprets internal variability as an absence of decisive policy news and therefore declines to revise its rate expectations.
